\documentclass[journal,10pt,doublecolumn]{IEEEtran}
\usepackage{cite}
\usepackage{amsmath,amssymb,amsfonts}
\usepackage{algorithmic}
\usepackage{graphicx}
\usepackage{textcomp}
\usepackage{xcolor}
\usepackage{booktabs}
\usepackage{multirow}
\usepackage{makecell}
\usepackage{hyperref}
\usepackage[caption=false,font=footnotesize]{subfig}

\begin{document}

\title{Restoration-Oriented Self-Supervised Priors via Partial Degradation MAE for Real-World Image Super-Resolution}

\author{Technical Report \\ 
University of Science and Technology of Hanoi (USTH) \\
Vietnam Academy of Science and Technology (VAST)}

\maketitle

\begin{abstract}
Real-world image super-resolution (SR) remains highly challenging due to complex, unknown degradations. Standard self-supervised methods like Masked Autoencoders (MAE) learn generic representations by reconstructing clean pixels from randomly masked inputs, which lack awareness of degradation patterns. To address this, we propose a novel framework, \textbf{Partial Degradation MAE (PD-MAE-SR)}, designed to learn restoration-oriented priors. Instead of blanking patches, PD-MAE degrades 75\% of complex regions using a realistic camera degradation pipeline while training the encoder to reconstruct the clean high-resolution (HR) target. To bridge the domain shift between synthetic degradations and clean textures, we introduce a two-stage pretraining curriculum: Stage 1 pretrains on partially degraded inputs, and Stage 2 calibrates the encoder on randomly masked clean HR inputs. In Stage 3, the pretrained PD-MAE encoder is frozen and integrated into a state-of-the-art Dense-Residual Connected Transformer (DRCT) backbone via Spatial Feature Transform (SFT) layers. Extensive evaluations at 500k iterations demonstrate that PD-MAE-DRCT significantly improves perceptual quality, achieving a NIQE reduction of \textbf{-0.85} on BSDS100 and \textbf{-0.38} on Urban100. Crucially, on real-world camera benchmarks (RealSR Canon/Nikon), our method breaks the conventional perception-distortion tradeoff, yielding simultaneous gains in both PSNR (+0.04/+0.05 dB) and NIQE (-0.31/-0.56) compared to the strong DRCT baseline.
\end{abstract}

\begin{IEEEkeywords}
Image Super-Resolution, Masked Autoencoder, Spatial Feature Transform, Real-World Degradation, Transformer.
\end{IEEEkeywords}

\section{Introduction}
Real-world image super-resolution (SR) aims to restore high-quality (HQ) details from low-quality (LQ) observations degraded by unknown mixtures of blur, noise, downsampling, and compression. While generative adversarial networks (GANs) and diffusion models have advanced the perceptual realism of restored images, they often introduce hallucinated artifacts or incur extreme computational costs during inference. On the other hand, feed-forward architectures like the Dense-Residual Connected Transformer (DRCT) exhibit high fidelity but struggle to generalize when faced with severe, out-of-distribution camera noise.

Recently, Self-Supervised Learning (SSL) via Masked Autoencoders (MAE) has emerged as a dominant paradigm in computer vision. However, standard MAE models are designed for downstream recognition tasks and suffer from two key limitations when applied to image restoration:
\begin{enumerate}
    \item \textbf{Degradation Unawareness:} Random patch masking removes content entirely, preventing the encoder from learning the subtle signatures of blur, sensor noise, or compression.
    \item \textbf{Domain Mismatch:} Reconstructing clean patches from masked clean inputs does not align with the blind SR objective, where the network must restore severely degraded regions.
\end{enumerate}

To overcome these challenges, we present \textbf{PD-MAE-SR}, a three-stage framework that repurposes MAE for blind super-resolution. The core contribution is a region-wise partial degradation pretraining scheme. By degrading selected high-frequency regions of a patch and forcing the model to reconstruct the intact HR counterpart, the encoder learns a robust restoration-oriented prior. We integrate this prior into the SOTA DRCT backbone using Spatial Feature Transform (SFT) layers. 

Our experimental results at 500,000 iterations validate this approach, showing significant perceptual improvements across multiple benchmarks and demonstrating a rare joint optimization of both PSNR and NIQE on real camera datasets.

\section{Proposed Method}
The training of the PD-MAE-SR framework is structured into three progressive stages, transitioning from self-supervised restoration pretraining to downstream backbone integration.

\subsection{Stage 1: Partial Degradation MAE Pretraining}
In the first stage, the objective is to train the encoder to recognize and correct realistic degradations. Unlike standard MAE which randomly masks patches, we employ a \textbf{Region-wise Partial Degradation} strategy based on edge complexity.

Let $I_{\text{HR}}$ be the high-resolution input patch of size $H \times W \times C$. The grayscale representation is denoted as $I_{\text{gray}}$. The horizontal and vertical gradients are computed using Sobel kernels:
\begin{equation}
G_x = I_{\text{gray}} * K_x, \quad G_y = I_{\text{gray}} * K_y
\end{equation}
where $*$ represents the 2D convolution operator, and $K_x, K_y$ are the horizontal and vertical Sobel operators, respectively. The spatial gradient magnitude $M$ is defined as:
\begin{equation}
M(i,j) = \sqrt{G_x(i,j)^2 + G_y(i,j)^2}
\end{equation}

Let $M_{\text{patch}}(k)$ be the average magnitude of patch index $k$. The patches are ranked based on $M_{\text{patch}}$, and the top $\theta = 75\%$ complex patches are selected for degradation:
\begin{equation}
\text{Mask}(k) = \begin{cases} 
   1 & \text{if } M_{\text{patch}}(k) \ge \tau \\ 
   0 & \text{otherwise} 
\end{cases}
\end{equation}
where $\tau$ is the $(1-\theta)$-th percentile threshold of gradient magnitudes across all patches in the image.

The selected patches are subjected to a realistic degradation pipeline (consisting of randomized blur, sensor noise, chromatic aberration, and JPEG compression). The remaining 25\% of patches are kept intact (clean HR).

The partially degraded image is fed into the MAE encoder-decoder. The network is optimized using an MSE loss calculated strictly on the degraded patches (Option X Loss), forcing the encoder to learn a restoration-oriented prior:
\begin{equation}
L_{\text{OptionX}} = \frac{1}{\sum_k \text{Mask}(k)} \sum_{k} \text{Mask}(k) \cdot \| \hat{I}_{\text{HR}}^{(k)} - I_{\text{HR}}^{(k)} \|_2^2
\end{equation}
where $\hat{I}_{\text{HR}}^{(k)}$ and $I_{\text{HR}}^{(k)}$ denote the reconstructed and ground-truth pixel values of the $k$-th patch, respectively.

\subsection{Stage 2: HR Structure Calibration}
Pretraining on heavily degraded synthetic inputs can cause the encoder to drift away from the clean HR texture domain. To calibrate the encoder, Stage 2 trains the network strictly on clean HR patches. We apply a 50\% random grid mask, setting the masked pixels to zero (black). The target remains the intact HR patch. The learning rate is reduced by a factor of 10 ($1 \times 10^{-5}$) to fine-tune the encoder weights gently, ensuring they remain representative of clean textures while retaining the restoration priors learned in Stage 1.

\subsection{Stage 3: Downstream SFT Injection into DRCT}
In the final stage, the pretrained and calibrated PD-MAE encoder is frozen and used as a structural guidance branch for the SR backbone (DRCT). 

The input LQ image is passed through the frozen PD-MAE encoder, producing a spatial feature map:
\begin{equation}
f_{\text{mae}} \in \mathbb{R}^{B \times 384 \times \frac{H}{p} \times \frac{W}{p}}
\end{equation}
where $p = 8$ is the patch size of the encoder. 

These features are projected to match the channel dimensions of DRCT and injected after key Dense-Residual Connected Transformer Groups (DCTG) using Spatial Feature Transform (SFT) layers:
\begin{equation}
\text{SFT}(F_{\text{drct}}, f_{\text{mae}}) = F_{\text{drct}} \odot (1 + \mathbf{S}(f_{\text{mae}})) + \mathbf{T}(f_{\text{mae}})
\end{equation}
where $\mathbf{S}$ and $\mathbf{T}$ represent the scale and shift mapping networks (implemented as $1 \times 1$ convolutional layers), and $\odot$ denotes element-wise multiplication.

To preserve baseline quality while guiding structure, training is guided by the backbone's native losses combined with an encoder consistency loss:
\begin{equation}
L_{\text{total}} = L_{\text{DRCT\_native}} + \lambda_{\text{mae}} L_{\text{PD\_MAE}}
\end{equation}
where $L_{\text{DRCT\_native}}$ is the original loss combination used by the DRCT backbone, $L_{\text{PD\_MAE}}$ is the consistency regularization term between intermediate feature layers of the backbone and the frozen MAE encoder, and $\lambda_{\text{mae}} = 0.05$ controls its influence. To avoid the catastrophic forgetting observed in initial experiments, Stage 3 utilizes the backbone's native online degradation pipeline.

\section{Experimental Results}
We trained both the baseline DRCT and the proposed PD-MAE-DRCT under identical environments for 500,000 iterations to perform a rigorous comparison.

\subsection{Quantitative Performance}
Table \ref{tab:results} summarizes the PSNR and NIQE metrics across various synthetic and real-world camera datasets.

\begin{table*}[htbp]
\centering
\caption{Quantitative Comparison at 500k Iterations (PSNR $\uparrow$ / NIQE $\downarrow$)}
\label{tab:results}
\begin{tabular}{lcccccc}
\toprule
\textbf{Model} & \textbf{Canon (RealSR)} & \textbf{Nikon (RealSR)} & \textbf{BSDS100} & \textbf{Urban100} & \textbf{Set14} & \textbf{Set5} \\
\midrule
DRCT Baseline & 27.18 / 7.14 & 26.59 / 7.34 & 23.93 / 7.42 & 21.59 / 6.01 & - / 7.37 & 27.59 / 7.85 \\
PD-MAE-DRCT (Ours) & \textbf{27.22} / \textbf{6.82} & \textbf{26.64} / \textbf{6.78} & 23.88 / \textbf{6.57} & 21.49 / \textbf{5.63} & - / \textbf{7.18} & 27.49 / 8.12 \\
\midrule
\textbf{Delta (Ours - Base)} & \textbf{+0.04} / \textbf{-0.31} & \textbf{+0.05} / \textbf{-0.56} & -0.05 / \textbf{-0.85} & -0.10 / \textbf{-0.38} & - / \textbf{-0.18} & -0.10 / +0.27 \\
\bottomrule
\end{tabular}
\end{table*}

\subsection{Qualitative Results}
To visually demonstrate the effectiveness of the PD-MAE prior, we present cropped regions of restored structures on real-world camera data in Fig. \ref{fig:qualitative_canon} and synthetic benchmarks in Fig. \ref{fig:qualitative_butterfly}. 

\begin{figure*}[!t]
    \centering
    \subfloat[Low Quality (LQ)]{\includegraphics[width=0.24\textwidth]{pd_mae/InputLR_Canon_004.png}\label{fig_lq_canon}}
    \hfil
    \subfloat[DRCT Baseline]{\includegraphics[width=0.24\textwidth]{pd_mae/Baseline_Canon_004_500000.png}\label{fig_baseline_canon}}
    \hfil
    \subfloat[PD-MAE-DRCT (Ours)]{\includegraphics[width=0.24\textwidth]{pd_mae/PD_MAE_Canon_004_500000.png}\label{fig_ours_canon}}
    \hfil
    \subfloat[Ground Truth (GT)]{\includegraphics[width=0.24\textwidth]{pd_mae/Canon_004.png}\label{fig_gt_canon}}
    \caption{Qualitative comparison on the RealSR (Canon\_004) dataset. The proposed PD-MAE-DRCT model successfully removes sensor noise and restores clean geometric structures, closely matching the Ground Truth reference without introducing ringing artifacts.}
    \label{fig:qualitative_canon}
\end{figure*}

\begin{figure*}[!t]
    \centering
    \subfloat[Low Quality (LQ)]{\includegraphics[width=0.24\textwidth]{pd_mae/InputLR_butterfly.png}\label{fig_lq_bf}}
    \hfil
    \subfloat[DRCT Baseline]{\includegraphics[width=0.24\textwidth]{pd_mae/Baseline_butterfly_500000.png}\label{fig_baseline_bf}}
    \hfil
    \subfloat[PD-MAE-DRCT (Ours)]{\includegraphics[width=0.24\textwidth]{pd_mae/PD_MAE_butterfly_500000.png}\label{fig_ours_bf}}
    \hfil
    \subfloat[Ground Truth (GT)]{\includegraphics[width=0.24\textwidth]{pd_mae/butterfly.png}\label{fig_gt_bf}}
    \caption{Qualitative comparison on the Set5 (butterfly) dataset. Our model restores sharper edges and clearer wing patterns, indicating that the restoration-oriented prior helps guide the SFT layers even under standard synthetic degradation scenarios.}
    \label{fig:qualitative_butterfly}
\end{figure*}

\begin{figure*}[!t]
    \centering
    \subfloat[Low Quality (LQ)]{\includegraphics[width=0.24\textwidth]{pd_mae/InputLR_Canon_014.png}\label{fig_lq_c14}}
    \hfil
    \subfloat[DRCT Baseline]{\includegraphics[width=0.24\textwidth]{pd_mae/Baseline_Canon_014_500000.png}\label{fig_baseline_c14}}
    \hfil
    \subfloat[PD-MAE-DRCT (Ours)]{\includegraphics[width=0.24\textwidth]{pd_mae/PD_MAE_Canon_014_500000.png}\label{fig_ours_c14}}
    \hfil
    \subfloat[Ground Truth (GT)]{\includegraphics[width=0.24\textwidth]{pd_mae/Canon_014.png}\label{fig_gt_c14}}
    \caption{Qualitative comparison on the RealSR (Canon\_014) dataset. PD-MAE-DRCT recovers high-frequency edge structures and suppresses chromatic noise artifacts more effectively than the DRCT baseline.}
    \label{fig:qualitative_canon014}
\end{figure*}

\begin{figure*}[!t]
    \centering
    \subfloat[Low Quality (LQ)]{\includegraphics[width=0.24\textwidth]{pd_mae/InputLR_Nikon_018.png}\label{fig_lq_nk18}}
    \hfil
    \subfloat[DRCT Baseline]{\includegraphics[width=0.24\textwidth]{pd_mae/Baseline_Nikon_018_500000.png}\label{fig_baseline_nk18}}
    \hfil
    \subfloat[PD-MAE-DRCT (Ours)]{\includegraphics[width=0.24\textwidth]{pd_mae/PD_MAE_Nikon_018_500000.png}\label{fig_ours_nk18}}
    \hfil
    \subfloat[Ground Truth (GT)]{\includegraphics[width=0.24\textwidth]{pd_mae/Nikon_018.png}\label{fig_gt_nk18}}
    \caption{Qualitative comparison on the RealSR (Nikon\_018) dataset. The PD-MAE encoder prior effectively guides the SFT layers to suppress Nikon sensor noise patterns while restoring fine structural details.}
    \label{fig:qualitative_nikon018}
\end{figure*}

\begin{figure*}[!t]
    \centering
    \subfloat[Low Quality (LQ)]{\includegraphics[width=0.24\textwidth]{pd_mae/InputLR_Nikon_050.png}\label{fig_lq_nk50}}
    \hfil
    \subfloat[DRCT Baseline]{\includegraphics[width=0.24\textwidth]{pd_mae/Baseline_Nikon_050_500000.png}\label{fig_baseline_nk50}}
    \hfil
    \subfloat[PD-MAE-DRCT (Ours)]{\includegraphics[width=0.24\textwidth]{pd_mae/PD_MAE_Nikon_050_500000.png}\label{fig_ours_nk50}}
    \hfil
    \subfloat[Ground Truth (GT)]{\includegraphics[width=0.24\textwidth]{pd_mae/Nikon_050.png}\label{fig_gt_nk50}}
    \caption{Qualitative comparison on the RealSR (Nikon\_050) dataset. Our method consistently produces sharper textures and cleaner backgrounds compared to the baseline under challenging real-world Nikon camera noise.}
    \label{fig:qualitative_nikon050}
\end{figure*}

\begin{figure*}[!t]
    \centering
    \subfloat[Low Quality (LQ)]{\includegraphics[width=0.24\textwidth]{pd_mae/InputLR_baby.png}\label{fig_lq_baby}}
    \hfil
    \subfloat[DRCT Baseline]{\includegraphics[width=0.24\textwidth]{pd_mae/Baseline_baby_500000.png}\label{fig_baseline_baby}}
    \hfil
    \subfloat[PD-MAE-DRCT (Ours)]{\includegraphics[width=0.24\textwidth]{pd_mae/PD_MAE_baby_500000.png}\label{fig_ours_baby}}
    \hfil
    \subfloat[Ground Truth (GT)]{\includegraphics[width=0.24\textwidth]{pd_mae/baby.png}\label{fig_gt_baby}}
    \caption{Qualitative comparison on the Set5 (baby) dataset. Despite the modest PSNR tradeoff on synthetic benchmarks, our model preserves facial skin texture and hair detail more naturally than the baseline.}
    \label{fig:qualitative_baby}
\end{figure*}

As shown in Figs. \ref{fig:qualitative_canon}--\ref{fig:qualitative_baby}, the baseline DRCT often suffers from residual blur or noise artifacts, especially on real-world camera data. By incorporating the restoration-oriented structural prior via SFT modulation, our model consistently reconstructs finer details, sharper edges, and cleaner textures that more closely align with the Ground Truth reference.


\subsection{Analysis and Discussion}
The experimental data reveals several critical behaviors of the proposed framework:
\begin{enumerate}
    \item \textbf{Significant Perceptual Gains:} PD-MAE-DRCT achieves a substantial reduction in NIQE scores across the major datasets. Notably, on BSDS100, the NIQE improves by \textbf{0.85} (an 11.5\% improvement), while on Urban100 it improves by \textbf{0.38}. This indicates that the features extracted by the frozen PD-MAE encoder provide superior structural and textual guidance, generating sharper edges and more realistic details.
    \item \textbf{Dual Improvements on Real Camera Data:} On the real-world Canon and Nikon datasets, the model improves \textit{both} fidelity (PSNR) and perceptual quality (NIQE). Specifically, Canon PSNR increases by \textbf{0.04 dB} alongside a \textbf{-0.31} NIQE reduction; Nikon PSNR increases by \textbf{0.05 dB} with a \textbf{-0.56} NIQE reduction. This indicates that the pretraining curriculum effectively bridges the sim-to-real gap, allowing the SFT layers to regularize the network against real camera noise.
    \item \textbf{Minimal Distortion Tradeoff:} On synthetic datasets (Set5, BSDS100, Urban100), the slight drop in PSNR (ranging from -0.05 dB to -0.10 dB) is extremely small and represents a highly favorable tradeoff given the massive improvements in visual naturalness (NIQE).
\end{enumerate}

\section{Conclusion}
This report presented the PD-MAE-SR framework, which utilizes region-wise partial degradation pretraining and HR calibration to extract restoration-oriented structural priors. By integrating these priors into the SOTA DRCT backbone via SFT modulation, we achieved substantial perceptual gains. The model demonstrates simultaneous improvements in PSNR and NIQE on real-world datasets, establishing a highly robust solution for blind image super-resolution without adding inference-time trainable parameters.

\end{document}
